% Hafta 4 — Kirli veri takibi: kaynaktan hedefe
% Derleme: py -3.12 tools/latex/derle.py docs/week-4/assets/h04-14-kirli-veri.tex
\documentclass[tikz,border=8pt]{standalone}
\usepackage{cen429-cizim}
\begin{document}
\begin{tikzpicture}[node distance=10mm and 14mm]
  \node[baslik] (b) at (0,0) {Kirli veri takibi: kaynaktan hedefe};
  \node[altbaslik] at (b.south west) {saldırganın yazdığı bir değer, denetlenmeden tehlikeli bir çağrıya ulaşıyor};

  \node[kotukutu=42mm, below=16mm of b.south west, anchor=north west] (kaynak) {%
    \kb{KAYNAK (source)}\\
    \kod{recv(sock, buf, ...)}\\
    ağdan gelen veri\\
    güvenilmez};
  \node[uyarikutu=42mm, right=of kaynak] (tasima) {%
    \kb{TAŞIMA}\\
    \kod{len = buf[1]<{}<8 | buf[2]}\\
    doğrulama YOK\\
    kirlilik korunur};
  \node[kotukutu=42mm, right=of tasima] (hedef) {%
    \kb{HEDEF (sink)}\\
    \kod{memcpy(out, buf+3, len)}\\
    sınır denetimi yok\\
    CWE-787};

  \draw[ok] (kaynak) -- (tasima);
  \draw[ok] (tasima) -- (hedef);

  \node[fit=(kaynak)(hedef), inner sep=0pt] (grup) {};
  \node[kotudip, text width=150mm, below=10mm of grup.south, anchor=north] {Bulgu kuralı: kaynak ile hedef arasında doğrulama yoksa, o yol bir zafiyettir.};
\end{tikzpicture}
\end{document}
